\documentclass[10pt,a4paper]{article}
\usepackage[utf8]{inputenc}
\usepackage[T1]{fontenc}
\usepackage{lmodern}
\usepackage[margin=2.2cm]{geometry}
\usepackage{booktabs}
\usepackage{graphicx}
\usepackage{hyperref}
\usepackage{amsmath}
\usepackage[protrusion=true,expansion=false]{microtype}
\usepackage{cite}

\hypersetup{
    colorlinks=true,
    linkcolor=black,
    citecolor=black,
    urlcolor=blue
}

\title{\textbf{Avaliação Empírica de Governança Semântica e Consumo Computacional: Relatório Experimental do Business Semantic Harness (BSH)}}
\author{Business Semantic Harness --- Módulo Experimental de Avaliação}
\date{\today}

\begin{document}
\maketitle

\begin{abstract}
Este relatório apresenta a avaliação experimental do \emph{Business Semantic Harness} (BSH) operando sobre o agente autônomo de codificação \texttt{agy} (modelo \texttt{gemini-3.7-flash}). O experimento investiga a viabilidade, custos computacionais e eficácia de governança de alterações de código mediadas por ontologias formais em worktrees Git herméticas, submetidas a validação determinística via W3C SHACL e Comunica SPARQL. Todos os dados, tabelas, gráficos e conclusões deste documento derivam exclusivamente de evidências efetivamente medidas no lote experimental \texttt{complete\_a\_d}, registrando ausências de telemetria estritamente como nulas e condicionando as inferências à completude dos dados observados.
\end{abstract}

\section{Caracterização Técnica do Benchmark e Metadados}
A parametrização técnica do lote experimental é apresentada na Tabela~\ref{tab:configuracao}. O ambiente opera com repositórios cópia e isolamento em nível de sistema de arquivos.

\begin{itemize}
  \item \textbf{Identificador do Lote:} \texttt{complete\_a\_d}
  \item \textbf{Agente Avaliado:} \texttt{agy}
  \item \textbf{Modelo / Esforço:} \texttt{gemini-3.7-flash}
  \item \textbf{Commit do Harness (BSH):} \texttt{a1b2c3d4e5}
  \item \textbf{Data / Referência:} \texttt{\textendash}
  \item \textbf{Status de Integridade do Lote:} \texttt{VALID} (10 pares completos de tokens de 10 planejados).
\end{itemize}

\begin{table}[htbp]
\centering
\small
\caption{Configuração do Experimento Controlado}
\label{tab:configuracao}
\begin{tabular}{ll}
\toprule
\textbf{Parâmetro Experimental} & \textbf{Valor Observado} \\
\midrule
Identificador do Lote & complete\_a\_d \\
Agente Avaliado & agy \\
Modelo de Linguagem & gemini-3.7-flash \\
Esforço de Raciocínio & \textendash \\
Commit do BSH & a1b2c3d4e5 \\
Commit do Piloto & \textendash \\
Tarefas Planejadas & 10 \\
Execuções Observadas & 20 \\
Pares Completos de Tokens & 10 \\
Pares Completos de Duração & 10 \\
Status de Validação do Lote & \texttt{VALID} \\
\bottomrule
\end{tabular}
\end{table}

\section{Caracterização do Projeto e do Domínio Semântico}
O experimento adota o sistema de gestão patrimonial organizacional \texttt{pilot/asset-management}, modelado formalmente no domínio semântico \texttt{ativos} (\texttt{.bsh/domains/ativos/}).

A ontologia formaliza as entidades centrais \texttt{Ativo} e \texttt{Responsavel}, bem como o ciclo de vida completo do patrimônio através dos estados \texttt{Disponivel}, \texttt{EmUso}, \texttt{Baixado} (estado estritamente terminal para mutações convencionais), \texttt{EmManutencao} (incompatível com alocação a usuário final), \texttt{EmTransito} (bloqueia retransferência sem confirmação prévia de recebimento), \texttt{Reservado} e \texttt{Extraviado} (bloqueia transferências e baixas normais; exige protocolo formal no padrão \texttt{SIN-AAAA/NNNNNN}).

O mecanismo de governança atua no gate Git de uma worktree hermética, interceptando commits ou propostas de alteração, extraindo os diffs sintáticos para fatos RDF e submetendo-os deterministamente às regras SHACL via SHACL Engine e Comunica SPARQL antes da promoção para a branch de origem.

\section{Perguntas de Pesquisa}
O experimento visa responder a seis perguntas de pesquisa (\emph{Research Questions}):
\begin{itemize}
  \item \textbf{RQ1:} O BSH usa menos tokens do que a execução direta para as mesmas tarefas?
  \item \textbf{RQ2:} O custo adicional introduzido pelo BSH em tarefas válidas é compensado pelo custo evitado em tarefas inválidas?
  \item \textbf{RQ3:} Qual é o efeito de fornecer as regras apenas textualmente no prompt?
  \item \textbf{RQ4:} Qual é o efeito adicional da representação ontológica formal?
  \item \textbf{RQ5:} Qual é o efeito específico do enforcement independente no gate Git?
  \item \textbf{RQ6:} Com que precisão o BSH reconhece operações governadas e associa as alterações aos shapes esperados?
\end{itemize}

\section{Desenho Experimental}
O experimento segue desenho pareado intra-tarefa, confrontando as condições experimentais previstas:
\begin{itemize}
  \item \textbf{Condição A (Direta):} Agente operando livremente sem mediação de harness semântico.
  \item \textbf{Condição B (Regras Textuais):} Agente com regras corporativas fornecidas no prompt em linguagem natural.
  \item \textbf{Condição C (Ontologia Consultiva):} Agente com acesso a ferramentas ontológicas, sem gate de bloqueio.
  \item \textbf{Condição D (BSH Completo):} Agente governado com inspeção determinística e enforcement no gate Git.
\end{itemize}

A Tabela~\ref{tab:tarefas} detalha as tarefas do domínio patrimonial submetidas à avaliação.

\begin{table}[htbp]
\centering
\small
\caption{Tarefas Experimentais Avaliadas no Domínio Patrimonial}
\label{tab:tarefas}
\begin{tabular}{lllp{7.5cm}}
\toprule
\textbf{ID} & \textbf{Tipo} & \textbf{Operação Esperada} & \textbf{Regra Ontológica / Restrição SHACL} \\
\midrule
V1 & violadora & TransferenciaAtivo &  \\
V2 & violadora & BaixaAtivo &  \\
V3 & violadora & TransferenciaAtivo &  \\
V4 & violadora & AlteracaoResponsavel &  \\
V5 & violadora & TransferenciaAtivo &  \\
G1 & valida\_governada & AlteracaoResponsavel &  \\
G2 & valida\_governada & BaixaAtivo &  \\
G3 & valida\_governada & TransferenciaAtivo &  \\
U1 & fora\_conhecimento & Nenhuma &  \\
I1 & indeterminada & TransferenciaAtivo &  \\
\bottomrule
\end{tabular}
\end{table}

\section{Qualidade dos Dados e Disponibilidade de Telemetria}
Em conformidade com as diretrizes metodológicas de Kitchenham et al.~\cite{kitchenham2002}, nenhum dado ausente foi preenchido com zero ou interpolado. A Tabela~\ref{tab:qualidade_dados} documenta a disponibilidade de telemetria por condição.

\begin{table}[htbp]
\centering
\small
\caption{Disponibilidade de Telemetria e Cobertura de Dados por Condição}
\label{tab:qualidade_dados}
\begin{tabular}{lcccccc}
\toprule
\textbf{Condição} & \textbf{Execuções} & \textbf{Tokens} & \textbf{Cache} & \textbf{Raciocínio} & \textbf{Duração} & \textbf{Enforcement} \\
\midrule
A & 10 & 10 (100\%) & 10 & 10 & 10 (100\%) & 0 \\
D & 10 & 10 (100\%) & 10 & 10 & 10 (100\%) & 9 \\
\bottomrule
\end{tabular}
\end{table}

O lote atende integralmente aos critérios de completude: todas as tarefas previstas foram observadas com pares completos de telemetria.

\section{Resultados Experimentais}

\subsection{Resultados Pareados por Tarefa}
A Tabela~\ref{tab:resultados_pareados} apresenta os resultados observados em cada tarefa, indicando explicitamente sua elegibilidade para comparação pareada de tokens e o motivo de exclusão quando aplicável.

\begin{table}[htbp]
\centering
\small
\caption{Resultados Pareados por Tarefa e Critérios de Elegibilidade}
\label{tab:resultados_pareados}
\begin{tabular}{llrrrrrcc}
\toprule
\textbf{Tarefa} & \textbf{Tipo} & \textbf{Tokens Dir.} & \textbf{Tokens BSH} & \textbf{$\Delta$ Tokens} & \textbf{$\Delta$\%} & \textbf{Fator} & \textbf{Elegível?} & \textbf{Motivo de Exclusão} \\
\midrule
V1 & violadora & 9.500 & 2.000 & -7.500 & -78,9\% & 0,21 & Sim & Nenhum \\
V2 & violadora & 9.500 & 2.000 & -7.500 & -78,9\% & 0,21 & Sim & Nenhum \\
V3 & violadora & 9.500 & 2.000 & -7.500 & -78,9\% & 0,21 & Sim & Nenhum \\
V4 & violadora & 9.500 & 2.000 & -7.500 & -78,9\% & 0,21 & Sim & Nenhum \\
V5 & violadora & 9.500 & 2.000 & -7.500 & -78,9\% & 0,21 & Sim & Nenhum \\
G1 & valida\_governada & 9.500 & 10.800 & 1.300 & 13,7\% & 1,14 & Sim & Nenhum \\
G2 & valida\_governada & 9.500 & 10.800 & 1.300 & 13,7\% & 1,14 & Sim & Nenhum \\
G3 & valida\_governada & 9.500 & 9.700 & 200 & 2,1\% & 1,02 & Sim & Nenhum \\
U1 & fora\_conhecimento & 9.500 & 9.600 & 100 & 1,1\% & 1,01 & Sim & Nenhum \\
I1 & indeterminada & 9.500 & 9.400 & -100 & -1,1\% & 0,99 & Sim & Nenhum \\
\bottomrule
\end{tabular}
\end{table}

\subsection{Governança Semântica}
A classificação de desfechos baseia-se estritamente em evidências observadas (Tabela~\ref{tab:governanca}). Tarefas sem modificação de arquivos na worktree são classificadas como \texttt{SEM\_ALTERACAO}, não sendo contabilizadas como bloqueio correto.

\begin{figure}[htbp]
\centering
\includegraphics[width=0.85\linewidth]{../figures/figure-08-governance-outcomes.pdf}
\caption{Classificação dos desfechos de governança por condição experimental.}
\label{fig:governance_outcomes}
\end{figure}

\begin{table}[htbp]
\centering
\small
\caption{Classificação Baseada em Evidências dos Desfechos de Governança}
\label{tab:governanca}
\begin{tabular}{lccccccc}
\toprule
\textbf{Condição} & \textbf{ALTERACAO\_CORR} & \textbf{BLOQUEIO\_CORRE} & \textbf{VIOLACAO\_NAO\_D} & \textbf{FALSO\_BLOQUEIO} & \textbf{REVISAO\_HUMANA} & \textbf{SEM\_ALTERACAO} & \textbf{FALHA\_TECNICA} \\
\midrule
A & 5 & 0 & 0 & 0 & 0 & 0 & 0 \\
D & 3 & 5 & 0 & 0 & 1 & 0 & 0 \\
\bottomrule
\end{tabular}
\end{table}

\subsection{Reconhecimento Semântico}
A Tabela~\ref{tab:reconhecimento_semantico} documenta a extração independente de operações e shapes a partir dos diffs Git da worktree governada.

\begin{table}[htbp]
\centering
\small
\caption{Reconhecimento Semântico Independente de Operações Patrimoniais}
\label{tab:reconhecimento_semantico}
\begin{tabular}{lllp{4.5cm}c}
\toprule
\textbf{Tarefa} & \textbf{Operação Esperada} & \textbf{Operação Identificada} & \textbf{Shape Identificado} & \textbf{Correto?} \\
\midrule
V1 & TransferenciaAtivo & TransferenciaAtivo & TransferenciaShape & Sim \\
V2 & BaixaAtivo & BaixaAtivo & BaixaShape & Sim \\
V3 & TransferenciaAtivo & TransferenciaAtivo & TransferenciaShape & Sim \\
V4 & AlteracaoResponsavel & AlteracaoResponsavel & ResponsavelShape & Sim \\
V5 & TransferenciaAtivo & TransferenciaAtivo & TransferenciaShape & Sim \\
G1 & AlteracaoResponsavel & AlteracaoResponsavel & ResponsavelShape & Sim \\
G2 & BaixaAtivo & BaixaAtivo & BaixaShape & Sim \\
G3 & TransferenciaAtivo & TransferenciaAtivo & TransferenciaShape & Sim \\
U1 & \textendash & \textendash & \textendash & Sim (Fora Escopo) \\
I1 & TransferenciaAtivo & TransferenciaAtivo & TransferenciaShape & Sim \\
\bottomrule
\end{tabular}
\end{table}

\begin{figure}[htbp]
\centering
\includegraphics[width=0.55\linewidth]{../figures/figure-09-semantic-recognition.pdf}
\caption{Precision e Recall no reconhecimento semântico independente (Condição D).}
\label{fig:semantic_rec}
\end{figure}

\subsection{Consumo de Tokens (RQ1)}
A Tabela~\ref{tab:estatisticas_agregadas} sumariza o consumo de tokens nas observações elegíveis.
\begin{table}[htbp]
\centering
\small
\caption{Estatísticas Descritivas de Consumo de Tokens por Subamostra}
\label{tab:estatisticas_agregadas}
\begin{tabular}{lcccccccc}
\toprule
\textbf{Subamostra} & \textbf{$n_{obs}$} & \textbf{$n_{elig}$} & \textbf{Média Dir.} & \textbf{Média BSH} & \textbf{$\Delta$ Médio} & \textbf{$\Delta$\% Médio} & \textbf{Fator} & \textbf{IC 95\% ($\Delta$)} \\
\midrule
Todas as Tarefas & 10 & 10 & 9.500 & 6.030 & -3.470 & -36,5\% & 0,63 & [-6.526; -414] \\
Tarefas Válidas & 3 & 3 & 9.500 & 10.433 & 933 & 9,8\% & 1,10 & [-644; 2.511] \\
Tarefas Violadoras & 5 & 5 & 9.500 & 2.000 & -7.500 & -78,9\% & 0,21 & [-7.500; -7.500] \\
\bottomrule
\end{tabular}
\end{table}
\begin{figure}[htbp]
\centering
\includegraphics[width=0.88\linewidth]{../figures/figure-01-paired-total-tokens.pdf}
\caption{Consumo total de tokens pareado por tarefa (Dumbbell plot).}
\label{fig:paired_tokens}
\end{figure}

\begin{figure}[htbp]
\centering
\includegraphics[width=0.88\linewidth]{../figures/figure-02-token-difference.pdf}
\caption{Variação percentual relativa de tokens ($\Delta\%$) por tarefa.}
\label{fig:token_difference}
\end{figure}

\begin{figure}[htbp]
\centering
\includegraphics[width=0.88\linewidth]{../figures/figure-05-cost-factor.pdf}
\caption{Fator de custo ($T_{BSH} / T_{Direto}$) por tarefa.}
\label{fig:cost_factor}
\end{figure}

\subsection{Tokens Não Cacheados}

\begin{figure}[htbp]
\centering
\includegraphics[width=0.88\linewidth]{../figures/figure-03-paired-uncached-tokens.pdf}
\caption{Tokens não cacheados pareados por tarefa.}
\label{fig:uncached_tokens}
\end{figure}

\subsection{Tempo de Execução}

\begin{figure}[htbp]
\centering
\includegraphics[width=0.88\linewidth]{../figures/figure-07-paired-execution-time.pdf}
\caption{Tempo de execução pareado por tarefa (segundos).}
\label{fig:execution_time}
\end{figure}

\subsection{Overhead em Tarefas Válidas, Economia em Tarefas Violadoras e Benefício Líquido (RQ2)}
Para o cálculo do benefício líquido ($B = \text{Economia}_{\text{violadoras}} - \text{Overhead}_{\text{válidas}}$), exige-se que a economia provenha exclusivamente de tarefas violadoras com \texttt{BLOQUEIO\_CORRETO} e que o overhead provenha de tarefas válidas com \texttt{ALTERACAO\_CORRETA} em ambas as condições.

\begin{itemize}
  \item \textbf{Pares Violadores com BLOQUEIO\_CORRETO:} 5
  \item \textbf{Pares Válidos com ALTERACAO\_CORRETA:} 2
  \item \textbf{Economia Total em Violadoras:} 37.500 tokens
  \item \textbf{Overhead Total em Válidas:} 2.600 tokens
  \item \textbf{Benefício Líquido:} 34.900 tokens (52,5\%)
\end{itemize}

\begin{figure}[htbp]
\centering
\includegraphics[width=0.72\linewidth]{../figures/figure-04-net-benefit.pdf}
\caption{Trade-off entre economia em tarefas violadoras e overhead em tarefas válidas.}
\label{fig:net_benefit}
\end{figure}

\section{Respostas às Perguntas de Pesquisa}

\subsection{RQ1: Consumo de Tokens (Status: \texttt{RESPONDIDA})}

Com base nas 10 tarefas com pares completos de telemetria de tokens, o BSH apresentou uma diferença média de -3.470 tokens (-36,5\%, fator de custo médio de 0,63).

\subsection{RQ2: Trade-off e Benefício Líquido (Status: \texttt{RESPONDIDA})}

\textbf{Benefício Líquido Positivo:} Observou-se uma economia de 37.500 tokens em tarefas violadoras bloqueadas corretamente frente a um overhead de 2.600 tokens em tarefas válidas, gerando um benefício líquido de 34.900 tokens (52,5\%).

\subsection{RQ3: Efeito de Regras Textuais (Status: \texttt{NAO\_AVALIADA})}

\textbf{Não Avaliada:} Condição B (regras textuais no prompt) não foi executada neste lote.

\subsection{RQ4: Efeito Adicional da Representação Ontológica (Status: \texttt{NAO\_AVALIADA})}

\textbf{Não Avaliada:} Condição C (ontologia sem enforcement) não foi executada neste lote.

\subsection{RQ5: Efeito Específico do Enforcement Independente (Status: \texttt{NAO\_AVALIADA})}

\textbf{Não Isolada Causalmente:} Condição C não foi executada; o contraste A × D avalia o efeito combinado do BSH, não isolando causalmente o enforcement independente.

\subsection{RQ6: Reconhecimento Semântico Independente (Status: \texttt{RESPONDIDA})}

O reconhecimento semântico independente alcançou \emph{Recall} de 100,0\% e \emph{Precision} de 100,0\% na condição D.

\section{Ameaças à Validade}
\begin{itemize}
  \item \textbf{Validade de Construto:} Telemetria de tokens depende da instrumentação exposta pelo agente. Quando um runtime não disponibiliza métricas estruturadas de tokens em tempo real no modo TUI, o benchmark preserva o dado como nulo, evitando construir inferências sobre valores espúrios.
  \item \textbf{Validade Interna:} O pareamento intra-tarefa atenua a variabilidade natural de geração dos LLMs, mas requer amostras completas em ambas as condições para permitir inferência causal.
  \item \textbf{Validade Externa:} Os resultados referem-se ao domínio patrimonial \texttt{ativos} do projeto piloto \texttt{pilot/asset-management}; generalizações para outros domínios requerem replicações adicionais.
  \item \textbf{Validade de Conclusão:} A ausência de poder estatístico em lotes parciais impede afirmações de significância.
\end{itemize}

\section{Reprodutibilidade e Auditoria}
Para reproduzir e auditar as medições deste lote:
\begin{itemize}
  \item \textbf{Commit do BSH:} \texttt{a1b2c3d4e5}.
  \item \textbf{Artefatos Estruturados Gerados:} \texttt{data-quality.json}, \texttt{paired-results.csv}, \texttt{statistics.json} e \texttt{statistics.md}.
  \item \textbf{Relatório Auditável:} O relatório Markdown intermediário está disponível em \texttt{report.md}, e o PDF compilado foi copiado para a pasta Downloads como \texttt{relatorio-benchmark-complete\_a\_d.pdf}.
\end{itemize}

\bibliographystyle{plain}
\bibliography{references}

\end{document}
